\documentclass[oneside,10pt]{report}
\usepackage[a4paper,left=2cm,right=2cm,top=2.5cm,bottom=2.5cm,footnotesep=0.5cm]{geometry}
\usepackage{ifthen}
\usepackage{verbatim}
\usepackage{graphpap}
\usepackage{graphicx}
\usepackage{pict2e}
\usepackage{color}
\usepackage{makeidx}
\usepackage{showidx}
\usepackage{latexsym}
\usepackage{amsmath}
\usepackage{multirow}
\usepackage{amssymb}
\usepackage{pdflscape}
\usepackage{listings}
\usepackage[dvipsnames]{xcolor}
\usepackage{caption}
\usepackage{subcaption}
\usepackage{tikz}
\usepackage{tabularx}
\usepackage{relsize}
\usepackage{enumitem}
\usepackage{float}
\usepackage{ragged2e}
\usepackage[colorlinks = true,
linkcolor = blue,
urlcolor  = blue,
citecolor = blue,
anchorcolor = blue]{hyperref}
\usepackage{url}
\usepackage{longtable}
%\usepackage{afterpage}
\usetikzlibrary{circuits.logic.US}
\usetikzlibrary{arrows.meta}
\allowdisplaybreaks[1]
\usepackage{pgfplots}
\usepackage{mathrsfs}
\usetikzlibrary{arrows}
\usetikzlibrary{shapes.geometric}
\usetikzlibrary{angles,quotes}
\pagestyle{empty}
\usepackage{array}
\usepackage{booktabs}
\usepackage{eso-pic}
\usepackage{pdfpages}

%\pgfplotsset{compat=1.15}
\newcommand{\degre}{\ensuremath{^\circ}}

%
% \mathset{set-name}
%     Stylized character corresponding to a math sets like 
% N - Natuaral Numbers
% Z - Integers
% Q - Rational Numbers
% R - Real Numbers
% C - Complex Numbers
%

\newcommand{\mathset}[1]
{%
    \mathbb
    {%
    {#1}%
    }%
}

\DeclareMathOperator{\arccot}{arccot}
\DeclareMathOperator{\arccsc}{arccsc}
\DeclareMathOperator{\arcsec}{arcsec}

\newcommand{\acTitle}[1]%
{%
    \chapter{#1}
}

\newcommand{\acHead}[1]%
{%
    \section{#1}
}

\newcommand{\union}
{%
    \cup
}

\newcommand{\bunion}
{%
    \bigcup
}

\newcommand{\intersect}
{%
    \cap
}

\newcommand{\bintersect}
{%
    \bigcap
}

\newcommand{\crossprod}
{%
    \times
}

\newcommand{\wherein}
{%
    \mid
}

\newcommand{\tendsto}
{%
    \rightarrow
}

\newcommand{\iorel}
{%
    \rightarrow
}

\newcommand{\convolves}
{%
    \star
}

\newcommand{\eqnlabel}[1]
{%
    \label{eq:#1}%
}

\newcommand{\eqnref}[1]
{%
    Eqn-(\ref{eq:#1})%
}

\newcommand{\axiomlabel}[1]
{%
    \label{axiom:#1}%
}

\newcommand{\axiomref}[1]
{%
    Axiom-(\ref{axiom:#1})%
}

\newcommand{\tablabel}[1]
{%
    \label{tab:#1}%
}

\newcommand{\tabref}[1]
{%
    Table-(\ref{tab:#1})%
}

\newcommand{\figlabel}[1]
{%
    \label{fig:#1}%
}

\newcommand{\figref}[1]
{%
    Figure-(\ref{fig:#1})%
}

\newcommand{\ptlabel}[1]
{%
    \label{pt:#1}%
}

\newcommand{\ptref}[1]
{%
    Point-(\ref{pt:#1})%
}

%
% LinkLabel
%     Label for location
%

\newcommand{\LinkLabel}[1]
{%
\hypertarget{#1}{ }
}

%
% LinkRef
%     Ref for section in the Act
%

\newcommand{\LinkRef}[2]
{%
    \hyperlink{#1}{#2}%
}


\newcommand{\incircle}[1]
{%
\begin{tikzpicture}
\node[draw,circle]{#1};
\end{tikzpicture}
}

\newcommand{\inrect}[1]
{%
\begin{tikzpicture}
\node[draw,rectangle]{#1};
\end{tikzpicture}
}

\newcommand{\indcircle}[1]
{%
\begin{tikzpicture}
\node[draw,circle,double]{#1};
\end{tikzpicture}
}

\newcommand{\drawbposs}
{%
\foreach \x in {0, 1, 2, 3, 4, 5, 6, 7, 8, 9, 10, 11, 12, 13, 14, 15}
    \draw (15 - \x + 0.5, 1) node[anchor=north] {\scriptsize \x};
}%

\newcommand{\xss}{0.6}
\newcommand{\yss}{0.2}

\newcommand{\drawbpos}
{%
\foreach \x in {0, 1, 2, 3, 4, 5, 6, 7, 8, 9, 10, 11, 12, 13, 14, 15, 16, 
               17, 18, 19, 20, 21, 22, 23, 24, 25, 26, 27, 28, 29, 30, 31}
    \draw (31 - \x + 0.5, 1) node[anchor=north] {\scriptsize \x};
}%

\newcommand{\xs}{0.325}
\newcommand{\ys}{0.18}

\newcommand{\euclidAxiomLabel}[1]
{%
    \label{euclidAxiomLabel:#1}%
}

\newcommand{\euclidAxiomRef}[1]
{%
    axiom #1 [\ref{euclidAxiomLabel:#1}]%
}

\newcommand{\euclidCmnNotionLabel}[1]
{%
    \label{euclidCommonNotionLabel:#1}%
}

\newcommand{\euclidCmnNotionRef}[1]
{%
    C.N #1 [\ref{euclidCommonNotionLabel:#1}]%
}

\newcommand{\euclidDefLabel}[2]
{%
    \label{euclidDefLabel:#1:#2}%
}

\newcommand{\euclidDefRef}[2]
{%
    definition #2 of book #1 (\ref{euclidDefLabel:#1:#2})%
}

\newcommand{\euclidPropLabel}[2]
{%
    \label{euclidPropLabel:#1:#2}%
}

\newcommand{\euclidPropRef}[2]
{%
    proposition #2 of book #1 (\ref{euclidPropLabel:#1:#2})%
}
\newcommand{\vect}[1]
{%
    \overrightarrow{#1}%
}

\newcommand{\nCr}[2]
{%
    ^{{#1}}C_{#2}%
}

\newcommand{\nPr}[2]
{%
    ^{#1}P_{#2}%
}

\tikzset{
    partial ellipse/.style args={#1:#2:#3}{
        insert path={+ (#1:#3) arc (#1:#2:#3)}
    }
}

\setlength{\parskip}{0.5\baselineskip}
\setlength{\parindent}{0pt}

\newboolean{newapi}
\setboolean{newapi}{false}
\newboolean{haveSolutions}

\newenvironment{apiprototype}%
{%
    \begin{lstlisting}[language=C]
}%
{%
    \end{lstlisting}
}

%
% \dsstart{component name}
%      Macro that is used for filling the contents in start of data structures
% section
%

\newcommand{\dsstart}[1]%
{%
This section provides information regarding various data structures and data 
buffers that application developer need to know for developing applications
using #1.
}

%
% \dsfields
%     Text added before data structure fields
%

\newcommand{\dsfields}%
{%
Below is brief description of fields in this structure.
}

%
% \apistart{component name}
%     Macro that is used for filling the contents in start of API section
%

\newcommand{\apistart}[1]
{%
This section provides information regarding different Application Program
Interfaces that is provided by #1.
}

%
% \apiparams{api name}
%     Text added before API parameter description
%

\newcommand{\apiparams}[1]
{%
Below is brief information on arguments of #1 function.
}

%
% \apirvalue{api name}
%     Text added before API return values
%

\newcommand{\apirvalue}[1]
{%
Possible return values from #1 are given below.
}

%
% \encapstruct{struct name}{component name}{category}
%     Macro that is used for filling the contents related to encapsulation
% structure in the data structures section
%

\newcommand{\encapstruct}[3]%
{%
\cstruct{#1} is created to encapsulate #2 implementation details to the 
application developer. Size of this structure is defined to be greater than or
equal to the corresponding internal structure used by the #3 implementation. As 
this structure is created only for encapsulation purpose, details of its field 
is not described here.
}

%--------------------------------------------------------------------------------------------------%
%                                                                                                  %
% Label and Reference commands:                                                                    %
%     This are commands that are used for creating a label and referencing them. They are there    %
% for following reasons                                                                            %
%    a. To have scope for different categories by adding a prefix for the same. In other words it  %
%       allows for \figlabel{topic1} and \tablabel{topic1}.                                        %
%    b. They use \mbox in the reference to ensure that the details and the reference number        %
%       doesn't appear in the same line.                                                           %
%                                                                                                  %
% Table below provides the list of them                                                            %
%                                                                                                  %
%     \seclabel, \secref: Sections                                                                 %
%     \eqlabel, \eqref: Equations                                                                  %
%     \tablabel, \tabref: For tables                                                               %
%     \figlabel, \figref: For figures                                                              %
%                                                                                                  %
%--------------------------------------------------------------------------------------------------%

%
% \seclabel{section-name}
%     Adds label for a section
%

\newcommand{\seclabel}[1]
{%
    \label
    {%
    {sec:#1}%
    }%
}

%
% \secref{section-name}
%     References label declared earlier
%

\newcommand{\secref}[1]
{%
    \mbox
    {%
    Section-\ref
        {%
        {sec:#1}%
        }%
    }%
}

%--------------------------------------------------------------------------------------------------%
%                                                                                                  %
% Audio Decoder/Parser User Manual usually have the following structure                            %
% Title                                                                                            %
%   |                                                                                              %
%   +---- Overview                                                                                 %
%   |     Subheading in these would have overview information about audio decoder/parser. They     %
%   |     need to use the macro                                                                    %
%   |         \prsrOverviewHead{heading}                                                           %
%   |                                                                                              %
%   +---- Data Structures and Buffers                                                              %
%   |     Subheadings in these would have information about individual data structure or buffer    %
%   |     information. They need to use the macro                                                  %
%   |         \adecDataStructHead{heading}                                                         %
%   |         \prsrDataStructHead{heading}                                                         %
%   |                                                                                              %
%   +---- Application Program Interface                                                            %
%   |     Subheading in these would have information about individual API related information.     %
%   |     They need to use the macro                                                               %
%   |         \adecApiHead{heading}                                                                %
%   |         \prsrApiHead{heading}                                                                %
%   |                                                                                              %
%   +---- Implementation Notes                                                                     %
%         Subheading in these would have information about specific point to note about the        %
%         implementation. They need to use the macro                                               % 
%             \adecInoteHead{heading}                                                              %
%             \prsrInoteHead{heading}                                                              %
%                                                                                                  %
% All the headings within the audio decoder/parser user manual need to go inside                   %
%     \adecHead,                                                                                   %
%     \prsrHead                                                                                    %
%                                                                                                  %
% Title pertaining to each audio decoder/parser would go inside \adecTitle, \prsrTitle             %
%                                                                                                  %
%--------------------------------------------------------------------------------------------------%

%
% \prsrOverviewHead{heading}
%     Heading in the parser overview section
%

\newcommand{\prsrOverviewHead}[1]
{%
    \subsubsection{#1}
}

%
% \adecDataStructHead{heading}
%     Audio Decoder - Data Structure heading
%

\newcommand{\adecDataStructHead}[1]
{%
    \subsubsection{#1}
}

%
% \prsrDataStructHead{heading}
%     Parser - Data Structure heading
%

\newcommand{\prsrDataStructHead}[1]
{%
    \subsubsection{#1}
}

%
% \adecApiHead{heading}
%     Audio Decoder - Application Program Interface heading
%

\newcommand{\adecApiHead}[1]
{%
    \subsubsection{#1}
}

%
% \prsrApiHead{heading}
%     Parser API heading
%

\newcommand{\prsrApiHead}[1]
{%
    \subsubsection{#1}
}

%
% \adecInoteHead{heading}
%     Audio Decoder - Implementation Notes heading
%

\newcommand{\adecInoteHead}[1]
{%
    \subsubsection{#1}
}

%
% \prsrInoteHead{heading}
%     Parser - Implementation Notes heading
%

\newcommand{\prsrInoteHead}[1]
{%
    \subsubsection{#1}
}

%
% \adecHead{heading}
%     Heading in the decoder user manual
%

\newcommand{\adecHead}[1]
{%
    \subsection{#1}
}

%
% \prsrHead{heading}
%     Heading in the parser user manual
%

\newcommand{\prsrHead}[1]
{%
    \subsection{#1}
}

%
% \adecTitle{title}
%     Decoder title level
%

\newcommand{\adecTitle}[1]
{%
    \section{#1}
}

%
% \prsrTitle{title}
%     Parser title level
%

\newcommand{\prsrTitle}[1]
{%
    \section{#1}
}

%
% \styleTitle{title}
%

\newcommand{\styleTitle}[1]
{%
    \section{#1}
}

%------------------------------------------------------------------------------

%
% \rectangle(basex,basey)(length,breadth)
%      Draw a rectangle
%

\def\rectangle(#1,#2)(#3,#4)%
{%
{%
    \put(#1,#2){\line(1,0){#3}}%
    \put(#1,#2){\line(0,1){#4}}%
    \count0=#1%
    \count1=#2%
    \advance\count0 by #3%
    \put(\count0,#2){\line(0,1){#4}}%
    \advance\count1 by #4%
    \put(\count0,\count1){\line(-1,0){#3}}%
}%
}

% Error code
\newcommand{\ecode}[1]{#1}

% C Language Macro
\newcommand{\cmacro}[1]{#1}

% C API
\newcommand{\capi}[1]{#1}

% C Data Structure
\newcommand{\cstruct}[1]{#1}

% Component Name
\newcommand{\cname}[1]{#1}

%
% datastruct
%     Environment used for providing information about data structures
%

\newenvironment{datastruct}
{%
\begin{center}
\begin{tabular}{|p{0.18\linewidth}|p{0.78\linewidth}|}
\hline
\multicolumn{1}{|c|}{\emph{Field}} & 
\multicolumn{1}{|c|}{\emph{Description}} \\ \hline
}%
{%
\end{tabular}%
\end{center}%
}

%
% describe
%     Environment for providing description
%

\newenvironment{describe}
{%
    \begin{list}{}{\setlength{\leftmargin}{0pt}}%
}%
{%
    \end{list}
}%

%
% \figfour{Input}{Box-1 Contents}{Box-2 Contents}
%         {Box-3 Contents}{Box-4 Contents}{Output}
%         {Caption}{Label}
% Draw a overview figure with four boxes
%

\newcommand{\figfour}[8]
{%
\begin{figure}[h]
\begin{center}
\setlength{\unitlength}{0.2em}
\begin{picture}(100,95)

% Input
\put(20,82){\makebox(60,10)[b]{\scriptsize #1}}

% Box-1 Contents
\rectangle(20,65)(60,10)%
\put(20,65){\makebox(60,10){\scriptsize #2}}

% Box-2 Contents
\rectangle(20,50)(60,10)%
\put(20,50){\makebox(60,10){\scriptsize #3}}

% Box-3 Contents
\rectangle(20,35)(60,10)
\put(20,35){\makebox(60,10){\scriptsize #4}}

% Box-4 Contents
\rectangle(20,20)(60,10)
\put(20,20){\makebox(60,10){\scriptsize #5}}

% Vectors
\put(50,20){\vector(0,-1){5}}
\put(50,80){\vector(0,-1){5}}
\put(50,65){\vector(0,-1){5}}
\put(50,50){\vector(0,-1){5}}
\put(50,35){\vector(0,-1){5}}

% Output
\put(20,3){\makebox(60,10)[t]{\scriptsize #6}}

\end{picture}
\end{center}
\caption{#7}
\label{#8}
\end{figure}
}

%
% \figfive{Input}{Box-1 Contents}{Box-2 Contents}
%         {Box-3 Contents}{Box-4 Contents}
%         {Box-5 Contents}{Output}
%         {Caption}{Label}
% Draw a overview figure with five boxes
%

\newcommand{\figfive}[9]
{%
\begin{figure}[h]
\begin{center}
\setlength{\unitlength}{0.2em}
\begin{picture}(100,100)
% Input
\put(20,91){\makebox(60,10)[b]{\scriptsize #1}}

% Box-1 Contents
\rectangle(20,75)(60,10)
\put(20,75){\makebox(60,10){\scriptsize #2}}

% Box-2 Contents
\rectangle(20,60)(60,10)%
\put(20,60){\makebox(60,10){\scriptsize #3}}

% Box-3 Contents
\rectangle(20,45)(60,10)
\put(20,45){\makebox(60,10){\scriptsize #4}}

% Box-4 Contents
\rectangle(20,30)(60,10)
\put(20,30){\makebox(60,10){\scriptsize #5}}

% Box-5 Contents
\rectangle(20,15)(60,10)
\put(20,15){\makebox(60,10){\scriptsize #6}}

% Output
\put(20,0){\makebox(60,9)[t]{\scriptsize #7}}

% Vectors
\put(50,90){\vector(0,-1){5}}
\put(50,75){\vector(0,-1){5}}
\put(50,60){\vector(0,-1){5}}
\put(50,45){\vector(0,-1){5}}
\put(50,30){\vector(0,-1){5}}
\put(50,15){\vector(0,-1){5}}
\end{picture}
\end{center}
\caption{#8}
\label{#9}
\end{figure}
}

\newcommand{\mc}[3]
{%
\multicolumn{#1}{#2}{#3}
}

\newcommand{\mr}[3]
{%
\multirow{#1}{#2}{#3}
}

\newcommand{\rou}[2]
{%
    \mathlarger{\mathlarger{\omega}}^{^{#1}}_{_{#2}}
}

\newcommand{\lconv}
{%
    \ast
}

\newcommand{\pconv}
{%
    \circledast_p
}

\newcommand{\cconv}
{%
    \circledast
}

\newcommand{\lxcorr}
{%
    \times
}

\newcommand{\pxcorr}
{%
    \otimes_p
}

\newcommand{\cxcorr}
{%
    \otimes
}

\newcommand{\Real}
{%
    \operatorname{Re}
}

\newcommand{\Imag}
{%
    \operatorname{Im}
}

\newcommand{\coss}[2]
{%
    C^{^{#1}}_{_{#2}}
}

\newcommand{\Sinc}
{%
    \operatorname{Sinc}
}

%
% levelone
%     Level information
%

\newcommand{\levelone}
{%
    \setlength{\leftskip}{0em}
}%

%
% leveltwo
%     Level two information
%

\newcommand{\leveltwo}
{%
    \setlength{\leftskip}{2em}
}

%
% levelthree
%     Level three of information
%

\newcommand{\levelthree}
{%
    \setlength{\leftskip}{4em}
}

%
% levelfour
%     Level four of information
%

\newcommand{\levelfour}
{%
    \setlength{\leftskip}{6em}
}

%
% SOPNote
%     SOP Notes
%

%
% SOPNoteHead
%     Heading for the SOP Notes
%

\newcommand{\SOPNoteHead}[1]
{%
    {%
    \setlength{\leftmargin}{0pt}
    \setlength{\parindent}{0em}
    \setlength{\leftskip}{0em}
    \par
    \textrm{%
    \textbf{%
    #1%
    }}}%
}

%
% SOPNoteHeadToc
%     Heading for the SOP Notes with Table of Contents
%

\newcommand{\SOPNoteHeadToc}[1]
{%
    {%
    \setlength{\leftmargin}{0pt}
    \setlength{\parindent}{0em}
    \setlength{\leftskip}{0em}
    \par
    \textrm{%
    \textbf{%
    #1%
    }}}%
    \addcontentsline{toc}{section}{#1}%
}

%
% SOPNoteBody
%     Body for the SOP Notes
%

\newcommand{\SOPNoteBody}
{%
    \mdseries
    \sffamily
    \setlength{\parskip}{1em}
    \setlength{\parindent}{0em}
    \setlength{\leftskip}{0em}
}

%
% SOPNoteFootNote
%     Section body for SOP Note
%

\newcommand{\SOPNoteFootNote}[1]
{%
\footnote{\sffamily\scriptsize #1}
}

%
% SOPNoteLabel
%     Label for SOP Note
%

\newcommand{\SOPNoteLabel}[1]
{%
\hypertarget{#1}{ }
}

%
% SOPNoteRef
%     Ref for SOP Note
%

\newcommand{\SOPNoteRef}[2]
{%
        \hyperlink{#1}{#2}%
}


%
% asr
%     Label for SOP Note
%

\newcommand{\asr}
{%
    \gg_a
}

\newcommand{\lsr}
{%
    \gg_l
}

\newcommand{\rsh}
{%
    \gg
}

\newcommand{\lsl}
{%
    \ll
}

\newcommand{\flags}[1]
{%
    \texttt{#1}%
}


\newcommand{\Lg}
{
    \log_2 N  % log base 2 of N
}                  
\newcommand{\NL}
{
    N\log_2 N  % N log2 N
}               
\newcommand{\hNL}
{
    \tfrac{N}{2}\log_2 N    % (N/2) log2 N  -- the per-butterfly unit
}    

\newcommand{\NLm}
{
    N(\log_2 N - 1)       % M85 butterfly-loop unit
}   

\newcommand{\E}{\;}                        % empty CM-85 cell filler


\newcommand{\ditto}{---''---}

\makeatletter
\newcommand\xleftrightarrow[2][]{%
\ext@arrow 9999{\longleftrightarrowfill@}{#1}{#2}}
\newcommand\longleftrightarrowfill@{%
\arrowfill@\leftarrow\relbar\rightarrow}
\makeatother

\tikzstyle{startstop} = [rectangle, rounded corners, 
minimum width=3cm, 
minimum height=1cm,
text centered, 
draw=black]

\tikzstyle{io} = [trapezium, 
trapezium stretches=true,
trapezium left angle=70, 
trapezium right angle=110, 
minimum width=3cm, 
minimum height=1cm, text centered, 
draw=black]

\tikzstyle{process} = [rectangle, 
minimum width=3cm, 
minimum height=1cm, 
text centered, 
text width=3cm, 
draw=black]

\tikzstyle{decision} = [diamond, 
minimum width=3cm, 
minimum height=1cm, 
text centered, 
draw=black]